\section{Statistics}
\label{sec:statistics}

Every section so far has computed with objects given exactly: an integer, a
polynomial, a function with a known derivative. Statistics begins at the other
end, with \emph{data}---a list of measured numbers, each carrying error, none of
them the answer on its own---and asks what can be said about the whole. The
questions are the classical ones. Where is the data centered? How far does it
spread? What is the shape of its tail? Do two quantities rise and fall together?
And how does a sample, once in hand, speak for the population it was drawn
from?\index{statistics!descriptive}

Mathilda answers these with the same principle that governs the rest of the
system: \emph{stay exact while the data is exact}. Hand \B{Variance} a list of
integers and it returns a rational number, not a rounded decimal; hand it a list
of measured reals and it returns a machine number, because that is all the data
supports. The estimators are the textbook ones---the sample variance carries the
unbiased \(n-1\) in its denominator, kurtosis is Pearson's (not the excess
form)---and every one threads columnwise over a matrix, so a table of several
variables is summarized in a single call. Underneath, when the data is a packed
array of machine reals, these reductions run on the buffer at the speed of C, and
each lowers inside \B{Compile}; the mathematics you type at the prompt is the
mathematics that runs in a tight numerical loop.

We take the subject in the order a data analyst would. First \emph{location}---the
mean and the median, and why the choice between them matters (\B{Mean},
\B{Median}, \B{Commonest}); then \emph{spread} (\B{Variance},
\B{StandardDeviation}, \B{RootMeanSquare}); the \emph{five-number summary} that
sketches a distribution with five order statistics (\B{Min}, \B{Quartiles},
\B{Max}); the \emph{shape} of the distribution through its moments (\B{Moment},
\B{CentralMoment}, \B{Skewness}, \B{Kurtosis}); the relationship between two
variables (\B{Covariance}, \B{Correlation}); the smoothing of a time series
(\B{MovingAverage}, \B{MovingMedian}, \B{ExponentialMovingAverage}); the tallying
of a sample into an empirical distribution (\B{Tally}, \B{Counts}); and we close
by crossing from description to \emph{probability}, drawing samples from a
distribution and measuring them back (\B{RandomVariate}, \B{PDF},
\B{LearnDistribution}).

\subsection{Location: the mean and the median}
\label{ssec:stat-location}

The oldest question in statistics is where a batch of numbers sits. The
\emph{mean} \(\bar x = \tfrac1n\sum x_i\) is the balance point; the \emph{median}
is the middle value once the data is sorted. On clean data they nearly agree, and
the choice between them looks like a matter of taste. It is not. Consider Simon
Newcomb's 1882 measurements of the time light takes to traverse a fixed
baseline\index{Newcomb, Simon}, recorded as deviations from \(24800\)
nanoseconds---and note the two readings, \(-44\) and \(-2\), that any modern eye
flags as blunders:

\begin{codepairs}
\pair{statistics/location/1}{Newcomb's first twenty timings, two of them wild.}
\pair{statistics/location/2}{\B{Mean} of integer data is exact: the rational
\(87/4\).}
\pair{statistics/location/3}{As a decimal, \(21.75\)---dragged down by the two
outliers.}
\pair{statistics/location/4}{\B{Median} sorts and takes the center: \(51/2 = 25.5\),
a full three units higher and untouched by the blunders.}
\pair{statistics/location/5}{\B{Commonest} reports the most frequent value---here a
tie between \(24\) and \(29\).}
\end{codepairs}

The gap between \(21.75\) and \(25.5\) is the whole lesson. The mean feels every
value in proportion to its size, so a single gross outlier moves it without limit;
the median feels only \emph{rank}, so up to half the data can be corrupted before
it budges. When a distribution is symmetric and clean the two coincide and the
mean is the more efficient estimate; when it is skewed or contaminated, the median
is the honest one.\index{median!robustness}

\calloutnote{Pitfall}{The mean is not robust. One fat-fingered data-entry error---a
misplaced decimal, a sign flip like Newcomb's \(-44\)---can shift a mean
arbitrarily far, while leaving the median essentially fixed. Report both when the
data might be dirty, and worry when they disagree.}

\subsection{Spread: variance, standard deviation, and RMS}
\label{ssec:stat-spread}

Location is half the story; the other half is how widely the data scatters about
it. The \emph{sample variance} averages the squared deviations from the mean, and
the \emph{standard deviation} is its square root, back in the original units.
Albert Michelson's 1879 speed-of-light experiment\index{Michelson, Albert} gives a
clean integer batch to work with---measurements in km/s above \(299000\):

\begin{codepairs}
\pair{statistics/spread/1}{Michelson's first twenty runs.}
\pair{statistics/spread/2}{Their mean, exact at \(909\) km/s above the baseline.}
\pair{statistics/spread/3}{\B{Variance} stays exact too: the rational
\(209180/19\).}
\pair{statistics/spread/4}{\B{StandardDeviation} is its exact square root, the surd
\(2\sqrt{52295/19}\).}
\pair{statistics/spread/5}{A decimal is one \B{N} away: about \(105\) km/s.}
\pair{statistics/spread/6}{\B{RootMeanSquare} is the quadratic mean
\(\sqrt{\tfrac1n\sum x_i^2}\)---magnitude about the origin, not about the mean.}
\end{codepairs}

Two details of the exact output repay attention. The variance is
\(209180/\mathbf{19}\), not over \(20\): the sample variance divides by \(n-1\),
not \(n\). And the standard deviation is returned as a genuine surd, carried
symbolically until you ask for digits---the same refusal to round prematurely that
runs through all of Mathilda.

\calloutnote{Theory}{Dividing by \(n-1\) rather than \(n\) is \emph{Bessel's
correction}\index{Bessel's correction}. The deviations are taken about the
\emph{sample} mean, which is itself pulled toward the data, so squared deviations
about it underestimate the true spread; the sample mean costs one degree of
freedom, and dividing by \(n-1\) restores an unbiased estimate of the population
variance.}

\usagebox{Variance}

\subsection{The five-number summary}
\label{ssec:stat-fivenum}

Before any model is fitted, a distribution can be sketched by five order
statistics: the minimum, the three quartiles, and the maximum. Together they
locate the center, bound the range, and measure the spread of the middle
half---the raw material of a box plot. Henry Cavendish's 1798 determination of the
mean density of the Earth\index{Cavendish, Henry} (relative to water), twenty-nine
measured values, is a natural specimen:

\begin{codepairs}
\pair{statistics/fivenum/1}{Cavendish's twenty-nine density measurements.}
\pair{statistics/fivenum/2}{\B{Min}: the lightest reading, \(4.88\).}
\pair{statistics/fivenum/3}{\B{Quartiles} returns \(\{Q_1, Q_2, Q_3\}\)---the
lower quartile, the median, and the upper quartile.}
\pair{statistics/fivenum/4}{\B{Max}: the heaviest, \(5.85\).}
\pair{statistics/fivenum/5}{The median again, \(5.46\), matching the middle
quartile.}
\pair{statistics/fivenum/6}{And the mean, \(5.448\)---almost equal to the median,
the signature of a roughly symmetric batch.}
\end{codepairs}

The interquartile range \(Q_3 - Q_1 = 5.6125 - 5.2975 = 0.315\) captures the
spread of the central half while ignoring the tails entirely, which is why it,
like the median, is a robust measure. That the mean and median nearly coincide
here---\(5.448\) against \(5.46\)---is the numerical fingerprint of symmetry;
where they part company, the data is skewed, and measuring that skew is the next
step.

\subsection{Shape: moments, skewness, and kurtosis}
\label{ssec:stat-shape}

Location and spread are the first two \emph{moments} of a distribution. The higher
moments describe its shape. The \(r\)-th raw moment is \(\tfrac1n\sum x_i^r\); the
\(r\)-th \emph{central} moment takes the deviations about the mean,
\(\tfrac1n\sum (x_i-\bar x)^r\). From these, two standardized quantities read off
the shape directly. Take a small, deliberately lopsided batch---the citation
counts of ten papers, nine ordinary and one runaway hit:

\begin{codepairs}
\pair{statistics/shape/1}{Ten citation counts, with one paper at \(20\).}
\pair{statistics/shape/2}{\B{Moment} gives the raw second moment
\(\tfrac1n\sum x_i^2 = 493/10\).}
\pair{statistics/shape/3}{\B{CentralMoment} of order \(2\)---the variance about the
mean, dividing by \(n\).}
\pair{statistics/shape/4}{The third central moment, \(43452/125\); its sign is the
sign of the skew.}
\end{codepairs}

\emph{Skewness} is the third central moment standardized by the cube of the
standard deviation; \emph{kurtosis} is the fourth, standardized by the square of
the variance. Both are pure numbers, free of units:

\begin{codepairs}
\pair{statistics/shape/5}{\B{Skewness}, exact---a positive surd, so the tail runs
to the right.}
\pair{statistics/shape/6}{About \(2.45\): strongly right-skewed, as the lone hit
demands.}
\pair{statistics/shape/7}{\B{Kurtosis}, the fourth standardized moment.}
\pair{statistics/shape/8}{About \(7.44\)---far above the \(3\) of a normal
distribution, the mark of a heavy tail.}
\end{codepairs}

A symmetric distribution has zero skewness; a positive value means a long right
tail, a negative value a long left one. Kurtosis measures how much of the variance
comes from rare, extreme deviations: Mathilda reports the \emph{Pearson} kurtosis,
for which the normal distribution sits at exactly \(3\), so the value \(7.44\) here
says the single outlier dominates the fourth moment.

\calloutnote{Theory}{Skewness and kurtosis are \emph{standardized} moments,
\(\mu_3/\sigma^3\) and \(\mu_4/\sigma^4\), built to be dimensionless and
invariant under shifting and rescaling the data. The ``excess kurtosis''
reported by some packages is simply Mathilda's kurtosis minus \(3\), centering
the normal distribution at zero.\index{kurtosis}\index{skewness}}

\subsection{Two variables: covariance and correlation}
\label{ssec:stat-correlation}

With two measurements on each subject, the question becomes whether they move
together. \B{Covariance} averages the product of paired deviations; \B{Correlation}
normalizes it by the two standard deviations to a pure number in \([-1, 1]\), so
that \(+1\) is perfect increasing linear agreement, \(-1\) perfect decreasing, and
\(0\) no linear relation at all.

No example makes the point---and its limits---better than \emph{Anscombe's
quartet}\index{Anscombe's quartet}, four data sets that Frank Anscombe constructed
in 1973 to be statistically indistinguishable by every summary yet obviously
different to the eye~\cite{anscombe1973}:

\begin{codepairs}
\pair{statistics/anscombe/1}{The common predictor \(x\) shared by the first three
sets.}
\pair{statistics/anscombe/2}{Set~I's response.}
\pair{statistics/anscombe/5}{Set~IV instead holds \(x\) nearly constant, with one
point far out.}
\pair{statistics/anscombe/6}{Set~IV's response.}
\end{codepairs}

Now compute the summaries of all four at once. They agree to the digits that
matter:

\begin{codepairs}
\pair{statistics/anscombe/7}{All four means of \(y\) sit at \(7.5\).}
\pair{statistics/anscombe/8}{All four variances of \(y\) sit at \(4.13\).}
\pair{statistics/anscombe/9}{And all four correlations equal \(0.816\).}
\pair{statistics/anscombe/10}{The covariance of set~I, \(5.501\)---unnormalized, in
the product of the units.}
\pair{statistics/anscombe/11}{Given the two columns as a matrix, \B{Correlation}
returns the correlation matrix: unit diagonal, \(0.816\) off it.}
\end{codepairs}

Yet the four data sets could hardly be more different: one is a clean line, one a
smooth curve, one a perfect line derailed by a single outlier, and one a vertical
stack whose slope is decided by one lone point. Identical means, identical
variances, identical correlations---and four unrelated stories. The moral, as old
as it is ignored, is that summary statistics never substitute for looking at the
data.

\calloutnote{Pitfall}{A correlation coefficient compresses a whole scatter of
points into one number, and much is lost in the compression. Anscombe's quartet is
the standing proof: equal correlations can hide a line, a curve, and an outlier.
Before trusting a coefficient, plot the points---\S\ref{sec:intro-stats} and the
graphics chapter show how.}

\usagebox{Correlation}

\subsection{Time series: moving statistics}
\label{ssec:stat-moving}

When the data is a sequence in time, its raw values jitter and the trend hides
underneath. A \emph{moving statistic} smooths the jitter by summarizing a sliding
window rather than the whole series at once. Take a short, noisy monthly series:

\begin{codepairs}
\pair{statistics/moving/1}{Twelve monthly readings, trending up through the noise.}
\pair{statistics/moving/2}{\B{MovingAverage} over a window of three---each output is
the mean of three neighbors, so the run is two shorter than the input.}
\pair{statistics/moving/3}{A \emph{weighted} moving average: weights \(\{1,2,1\}\)
lean on the center of each window.}
\pair{statistics/moving/4}{\B{MovingMedian} smooths the same way but resists a lone
spike, exactly as the median did in \S\ref{ssec:stat-location}.}
\pair{statistics/moving/5}{\B{ExponentialMovingAverage} weights every past value,
decaying geometrically; with \(\alpha=1/3\) the result stays exact.}
\pair{statistics/moving/6}{The same recurrence numerically, with \(\alpha=0.2\):
each output is \(0.2\) of the new value plus \(0.8\) of the running average.}
\end{codepairs}

The three smoothers trade off differently. The plain moving average is simplest
but lags the trend and reacts sharply as a window enters or leaves an outlier. The
moving median inherits the median's robustness, flattening isolated spikes
outright. The exponential moving average never forgets---every past point
contributes, weighted by \((1-\alpha)^k\)---so it responds smoothly and needs only
the previous output to advance, which is why it is the workhorse of streaming and
financial data.

\calloutnote{Performance}{These reductions are exactly the kind of numerical
loop Mathilda accelerates. On a packed array of machine reals the moving
statistics run on the raw buffer without boxing a single element, and the same
kernels lower inside \B{Compile}, so a smoothing pass over a million-point series
costs a single sweep of memory rather than a million interpreter steps.}

\subsection{Frequencies and the empirical distribution}
\label{ssec:stat-frequencies}

For discrete or categorical data the natural summary is not a mean but a
\emph{count}: how often each distinct value occurs. This tabulation is the
\emph{empirical distribution}, the sample's own stand-in for a probability
distribution. Fifteen rolls of a die will serve:

\begin{codepairs}
\pair{statistics/frequencies/1}{Fifteen die rolls.}
\pair{statistics/frequencies/2}{\B{Tally} pairs each distinct value with its count,
in order of first appearance.}
\pair{statistics/frequencies/3}{\B{Counts} returns the same tabulation as an
association, keyed by value---ready for \(O(1)\) lookup.}
\pair{statistics/frequencies/4}{\B{Commonest} extracts the mode, the most frequent
value.}
\pair{statistics/frequencies/5}{Asked for two, it returns the top two by
frequency.}
\end{codepairs}

Normalizing the counts by the sample size turns the tally into an empirical
probability mass function, and its running sum into an empirical cumulative
distribution---the discrete companions of the density and distribution functions
of the next subsection. For continuous data the same idea, applied to bins rather
than exact values, is a histogram, which the graphics chapter draws directly with
\B{Histogram}.

\subsection{From data to distributions}
\label{ssec:stat-distributions}

Descriptive statistics run in one direction, from a sample to its summaries.
Probability runs in the other, from a model to the samples it generates---and the
two meet in the estimators of this section, which are exactly the quantities that
converge to a distribution's parameters as the sample grows. Mathilda lets you
watch that convergence. Draw a thousand values from a normal population with mean
\(100\) and standard deviation \(15\)---the scale of an IQ test---and measure the
sample back:

\begin{codepairs}
\pair{statistics/distributions/1}{\B{SeedRandom} fixes the stream, so the draw is
reproducible from run to run.}
\pair{statistics/distributions/2}{A thousand draws from \B{NormalDistribution};
the semicolon holds back the long list.}
\pair{statistics/distributions/3}{The sample mean, \(100.20\)---within a fifth of a
unit of the true \(100\).}
\pair{statistics/distributions/4}{The sample standard deviation, \(15.03\), all but
recovering the true \(15\).}
\pair{statistics/distributions/5}{\B{PDF} evaluates the density; it is symmetric
about the mean, so \(85\) and \(115\) share a value and \(100\) peaks.}
\pair{statistics/distributions/6}{\B{LearnDistribution} runs the inference the
other way, fitting a distribution object back to the data.}
\end{codepairs}

That the sample mean and standard deviation land so near \(100\) and \(15\) is the
\emph{law of large numbers}\index{law of large numbers} in miniature: the sample
estimators converge to the population parameters as \(n\) grows, and a thousand
draws already pin them to a fraction of a unit. \B{RandomVariate} draws from a
distribution, \B{PDF} evaluates its density, and \B{LearnDistribution} inverts the
arrow---inferring a distribution from data---closing the loop between the
descriptive statistics of this section and the probability that underwrites them.

\calloutnote{Under the hood}{\B{RandomVariate} shares its stream with
\B{RandomReal}, so a single \B{SeedRandom} makes an entire session's simulation
reproducible---indispensable when a result must be regenerated exactly, as every
example in this book is. The normal draws come from a standard transform of that
uniform stream.}

\subsection*{Where this connects}

Statistics is where Mathilda's numerical substrate meets its data. Its estimators
lean directly on the reductions of the array engine---a \B{Mean} or \B{Variance}
of a packed column is a single buffer sweep, and each lowers inside \B{Compile} so
it runs at full speed inside a numerical loop (Chapter~\ref{ch:compilation}). The
correlation matrix of \S\ref{ssec:stat-correlation} is a Gram matrix, computed
through the same \B{Dot} and \B{Transpose} that drive the linear algebra of
\S\ref{sec:linalg}, and it is the first step of principal component analysis. The
sampling of \S\ref{ssec:stat-distributions} rests on the random-number generators,
and the fit performed by \B{LearnDistribution} opens into the machine-learning
family---classifiers, clustering, and density estimation---that treats a data set
not as something to summarize but as something to learn from. For the exhaustive
options of any function named here, its reference page is one \texttt{?Name} away
at the prompt.
